% Generated draft: requires mathematical and source review.
\documentclass[11pt]{article}
\usepackage[margin=25mm]{geometry}
\usepackage{fontspec}
\setmainfont{Latin Modern Roman}
\usepackage{amsmath,amssymb,mathtools,unicode-math}
\setmathfont{Latin Modern Math}
\usepackage{enumitem,needspace,xurl,hyperref}
\setlength{\parindent}{0pt}
\setlength{\parskip}{5pt}
\setlength{\emergencystretch}{3em}
\newcommand{\N}{\mathbb{N}}
\newcommand{\Z}{\mathbb{Z}}
\newcommand{\Q}{\mathbb{Q}}
\newcommand{\R}{\mathbb{R}}
\newcommand{\C}{\mathbb{C}}
\newcommand{\F}{\mathbb{F}}
\newcommand{\A}{\mathbb{A}}
\newcommand{\PP}{\mathbb{P}}
\newcommand{\E}{\mathbb{E}}
\newcommand{\eps}{\varepsilon}
\newcommand{\dd}{\,\mathrm{d}}
\begin{document}
\section*{UnsolvedMath ID: 30004840}
\textbf{Source problem code:} catalogue entry
\subsection{Definitions and mathematical statement}
For a compact quantum group \(\mathbb G\), write \(A=L^\infty(\mathbb G)\) and let \(\Delta:A\to A\,\bar\otimes\,A\) denote its coproduct, where \(\bar\otimes\) is the spatial von Neumann algebra tensor product. To express the questions, use the right-coaction convention: a coaction on a von Neumann algebra \(M\) is a normal injective unital \(*\)-homomorphism \(\alpha:M\to M\,\bar\otimes\,A\) satisfying
\[
(\alpha\otimes\mathrm{id})\alpha=(\mathrm{id}\otimes\Delta)\alpha.
\]
It is ergodic if
\[
M^\alpha:=\{x\in M:\alpha(x)=x\otimes1\}=\C1.
\]
In this convention a coideal von Neumann subalgebra is a unital von Neumann subalgebra \(N\subseteq A\) such that \(\Delta(N)\subseteq N\,\bar\otimes\,A\), equipped with the restricted coproduct. The supplied record does not specify the coaction side; this convention makes the notation explicit and does not reconstruct the report's conventions.

Two coactions of the same quantum group are equivariantly \(W^*\)-Morita equivalent if they can be realized as the two diagonal corners of a linking von Neumann algebra carrying a coaction: the complementary corner projections are fixed by the coaction, each has central support \(1\), and the corner identifications intertwine the coactions. Denote this relation by \(\sim_{\mathbb G}\). A type-I von Neumann algebra is one in which every nonzero central summand contains a nonzero abelian projection. A type-I factor is isomorphic to \(B(H)\) for a nonzero Hilbert space \(H\); the supplied record imposes no dimension or separability restriction on \(H\).

\textbf{Question in the source's parameter notation.} Write \(A_F=L^\infty(FO_F)\). For every ergodic coaction \((M_F,\alpha_F)\) covered by the source's parameter conventions, is there a parameter \(G\in F_q\) and a coideal von Neumann subalgebra \(N\subseteq A_G\) such that
\[
(M_G,\alpha_G)\sim_{FO_G}(N,\Delta_G|_N)?
\]
Here \((M_G,\alpha_G)\) must be the coaction associated with the given \((M_F,\alpha_F)\) by the source's intended construction. It is not an independently chosen coaction. The supplied record does not define \(F_q\), the admissible parameters \(F\), the quantum groups \(FO_F\), or this association. Consequently this question remains a conditional formulation in the source's notation, rather than a fully specified assertion. The cleaned statement additionally asks, if the universal assertion fails, for a characterization of the coactions having such a coideal representative; that follow-up is not present in the supplied original excerpt.

\textbf{Question for quantum \(SU_q(2)\).} Put \(A_q=L^\infty(SU_q(2))\). For every ergodic coaction \(\alpha:M\to M\,\bar\otimes\,A_q\) on a type-I von Neumann algebra \(M\), does there exist a coideal von Neumann subalgebra \(N\subseteq A_q\) such that
\[
(M,\alpha)\sim_{SU_q(2)}(N,\Delta_q|_N)?
\]
The admissible values of \(q\) and the convention for \(SU_q(2)\) are not supplied, so the quantifier over \(q\) cannot be completed from this record.

\textbf{Scope qualification.} The original excerpt motivates the first question by classification of projective unitary representations of quantum \(SU(2)\), described there as coactions of \(L^\infty(SU_q(2))\) on type-I factors \(B(H)\). It attributes explicit knowledge of the von Neumann algebra structure of coideals to [Tom08] and the projective-representation interpretation to [DC11]. The cleaned question instead explicitly assumes ergodicity and permits any type-I von Neumann algebra. These targets are not identified here: the supplied excerpt does not establish their equivalence or supply the missing parameter-transport construction. Its assertion that the first question would yield the stated classification is retained as source motivation, not proved as part of this draft.
\subsection{Short English statement}
The questions ask whether the specified ergodic quantum-group coactions have representatives, up to equivariant \(W^*\)-Morita equivalence, that arise from coideal von Neumann subalgebras. The general \(FO_F\) question lacks definitions of its parameter family and of the passage from \(M_F\) to \(M_G\). The \(SU_q(2)\) question lacks a range for \(q\), and its type-I-algebra formulation differs from the original excerpt's projective-representation motivation on factors \(B(H)\). These qualifications require mathematical review.
\subsection{Sources}
\begin{itemize}
\item Local UnsolvedMath registry, \texttt{data/unsolvedmath/problems.json}, permanent ID \(30004840\), source problem code \texttt{catalogue entry}, title ``Morita Classification of Quantum SU-Two Projective Representations.'' The supplied record includes both an original excerpt and a cleaned formulation.
\item Supplied primary-source citation: \emph{Quantum Groups -- Algebra, Analysis and Category Theory} (2022), Oberwolfach Reports, pp.\ 2397--2458, DOI \(10.4171/owr/2021/44\), \url{https://doi.org/10.4171/owr/2021/44}. The supplied excerpt carries page number 2451. The registry gives proposed year 2021 and names no proposer; its report citation gives 2022.
\item Supplied literature source, ``Oberwolfach Report 44'': \url{https://doi.org/10.4171/owr/2021/44}.
\item Supplied literature source, ``Projective representations for compact quantum groups'': \url{https://arxiv.org/abs/1006.2278}.
\item Supplied literature source, ``Projective representation theory and Baum--Connes'': \url{https://arxiv.org/abs/2112.04365}.
\item The original excerpt cites [Tom08] for explicit coideal von Neumann algebra structures and [DC11] for the interpretation of projective unitary representations as coactions on \(B(H)\). Full bibliographic entries and a verified correspondence with the supplied links are absent.
\item Literature assessment dated exactly \texttt{2026-08-22}: the supplied curation labels the problem ``partially solved'' and states that projective representations are characterized through actions on type-I factors and cocycle regularity, with explicit \(SU_q(2)\) examples, while a complete classification of all such coactions remains unavailable. This is supplied dated curation derived from the OWR collection, not independently verified mathematical or current-status information. The original excerpt describes the classification problem as open at its source date.
\item The registry labels the statement \texttt{corrected\_verified} and reports that the linked report was inspected and both bullet-point questions recovered with their hypotheses. This is a supplied verification claim; no source visit or independent verification is asserted in this draft.
\end{itemize}
\end{document}
